\begin{solapartat}
\punts{0,4} Sabem que les línies de camp surten de les càrregues positives i van a morir a les càrregues negatives.

\punts{0,15} Per tant, l'ió positiu es trobarà on hi hagi una font de línies de camp. A partir del gràfic i aproximadament, podem situar una font de línies de camp al punt $(-1,0)$~m

\punts{0,15} En canvi, l'ió negatiu es trobarà on hi hagi un pou de línies de camp. A partir del gràfic i aproximadament, podem situar un pou de línies de camp al punt $(1,0)$~m.

\punts{0,35} El nombre de línies de camp que surt d'una càrrega puntual (o va a morir a una càrrega puntual) és proporcional al valor absolut de la càrrega.

\punts{0,1} Podem comprovar que surten més línies de camp de la posició on es troba l'ió carregat positivament, $(-1,0)$, que de la posició on es troba l'ió carregat negativament, $(1,0)$.

\punts{0,1} per tant a $(-1,0)$ tenim una càrrega $+4e$ i a $(1,0)$ tenim una càrrega $-e$.
\end{solapartat}

\begin{solapartat}
\punts{0,65} Si ens situem a una distància molt més gran que al separació de les dues càrregues, el que veuríem es una càrrega puntual de valor igual a la suma de les dues càrregues, és a dir, veuríem el camp creat per una càrrega puntual de magnitud $+3e$.

\punts{0,6} Per tant, la superfície equipotencial s'aproximarà al d'una càrrega puntual, és a dir, serà una \destaca{superfície esfèrica}.
\end{solapartat}
